%%%
%
% $Autor: Siddhanth Sharma $
% $Datum: 2026-01-02 $
%
% !TeX spellcheck = GB
% !TeX program = pdflatex
% !TeX encoding = utf8
%
%%%

\chapter{Comparative Analysis: Edge Impulse vs.\ Our Implementation}
\label{sec:edge_impulse_compare}
\index{Edge Impulse}
\index{Deployment!Comparison}

\noindent
This chapter compares a \textbf{code-first} license plate recognition (LPR) implementation on an NVIDIA Jetson Nano with a \textbf{platform-first} workflow using \textit{Edge Impulse}.
The comparison is conducted at the level of the \textbf{end-to-end workflow}, not only model accuracy:
data acquisition and labeling, training and evaluation, deployment packaging, on-device profiling, and the cost of iteration.

\noindent
Edge Impulse is a web-based platform (Studio) for building, training, and deploying edge ML models, complemented by CLI tools and SDKs \cite{EdgeImpulseStudioOverview}.
It provides a structured ``Impulse'' concept that captures the ML pipeline as input + processing + learning blocks and packages this pipeline into deployment artefacts \cite{EdgeImpulseImpulseDesign}.
For our use case, the most important insight is that \textbf{Edge Impulse primarily standardizes the model lifecycle}, while our project primarily standardizes the \textbf{runtime application} (OCR heuristics, validation logic, billing state, dashboard integration).

\begin{figure}[htbp]
	\centering
	% in preamble (already mostly there)
% \usetikzlibrary{matrix,positioning}

\begin{tikzpicture}[
	font=\small,
	box/.style={draw, rounded corners, align=center, fill=blue!8,
		text width=42mm, inner sep=4pt},
	box2/.style={draw, rounded corners, align=center, fill=green!10,
		text width=42mm, inner sep=4pt},
	note/.style={draw, rounded corners, align=left, fill=gray!10,
		text width=110mm, inner sep=5pt},
	arrow/.style={-latex, thick},
	coltitle/.style={font=\bfseries, align=center}
	]
	
	% --- Column titles (anchored to columns)
	\node[coltitle] (t1) at (-4.6, 2.2) {Our Jetson LPR implementation\\(code-first)};
	\node[coltitle] (t2) at ( 4.6, 2.2) {Edge Impulse approach\\(platform-first)};
	
	% --- Left column (relative positioning, consistent spacing)
	\node[box, below=6mm of t1] (cam)  {Camera frames\\(\FILE{/dev/video0})};
	\node[box, below=5mm of cam] (det) {Plate detection\\YOLO (\mbox{Ultralytics} / TensorRT)};
	\node[box, below=5mm of det] (crop) {Crop \& preprocess\\(OpenCV, CLAHE, thresholding)};
	\node[box, below=5mm of crop] (ocr) {OCR\\EasyOCR (CPU/GPU)};
	\node[box, below=5mm of ocr] (logic) {Validation \& business logic\\(DE format, cooldown, billing)};
	\node[box, below=5mm of logic] (ui) {Storage \& UI\\SQLite + Dashboard\\(Flet / Flask)};
	
	\draw[arrow] (cam) -- (det);
	\draw[arrow] (det) -- (crop);
	\draw[arrow] (crop) -- (ocr);
	\draw[arrow] (ocr) -- (logic);
	\draw[arrow] (logic) -- (ui);
	
	% --- Right column aligned to left column (same y positions)
	\node[box2, right=26mm of cam] (da)
	{Data acquisition\\(Studio / CLI / API)};
	\node[box2, below=5mm of da] (lab)
	{Labeling\\(bounding boxes, AI-assisted)};
	\node[box2, below=5mm of lab] (imp)
	{Impulse design\\Input + Processing\\+ Learning blocks};
	\node[box2, below=5mm of imp] (train)
	{Train \& evaluate\\(built-in blocks or\\custom / BYOM)};
	\node[box2, below=5mm of train] (deploy)
	{Deploy artefact\\(Linux runner / EIM / SDK)};
	\node[box2, below=5mm of deploy] (app)
	{Application glue\\OCR + billing + UI\\remain custom};
	
	\draw[arrow] (da) -- (lab);
	\draw[arrow] (lab) -- (imp);
	\draw[arrow] (imp) -- (train);
	\draw[arrow] (train) -- (deploy);
	\draw[arrow] (deploy) -- (app);
	
	% --- Note: placed relative to the lowest nodes (no overlap)
	\node[note, below=10mm of $(ui.south)!0.5!(app.south)$] (note1) {%
		\textbf{Key comparison idea:}
		Edge Impulse standardizes the \emph{model lifecycle} (data, labeling, training, deployment packaging),
		while our project standardizes the \emph{application runtime} (OCR heuristics, validation rules,
		billing state, UI integration).};
	
\end{tikzpicture}
	\caption{Conceptual pipeline comparison: our Jetson LPR implementation vs.\ an Edge Impulse based workflow.}
	\label{fig:edge_impulse_pipeline_compare}
\end{figure}

% --------------------------------------------------------
\section{Scope and assumptions}
\label{sec:edge_impulse_scope}
\index{Scope}

\noindent
To keep the comparison concrete, we scope the study to:
\begin{itemize}
	\item \textbf{Task:} detect German license plates in camera frames and output a normalized plate string (plus application-side events such as billing entries).
	\item \textbf{Device class:} Linux microprocessor edge device (Jetson Nano), not a microcontroller.
	\item \textbf{Pipeline decomposition:} the LPR system is explicitly a \textbf{multi-stage pipeline}:
	(1) plate detection (object detection), (2) plate crop preprocessing, (3) OCR, (4) validation and business logic.
\end{itemize}

\noindent
This decomposition matters because Edge Impulse's built-in workflows naturally cover (1) and parts of preprocessing, but \textbf{OCR and domain validation remain application logic} in most practical deployments.
Therefore, a fair comparison focuses on \textbf{what changes} if the detection model lifecycle is moved to Edge Impulse, while OCR and downstream logic remain in our own code.

% --------------------------------------------------------
\section{Baseline: our implementation (what exists in this repository)}
\label{sec:edge_impulse_baseline}
\index{Baseline}

\noindent
Our project is implemented as a Python-first system optimized for Jetson Nano constraints:
\begin{itemize}
	\item \textbf{Detection model training:} YOLO-based detector trained on a Roboflow German license plate dataset.\\
	\noindent See project overview: \PATH{Code/LicencePlateDetection/about.md}.
	
	\item \textbf{Jetson runtime variant A (Ultralytics):} Live detection with Ultralytics YOLO.\\
	\noindent Implementation: \PATH{Code/LicencePlateDetection/jetson\_app/detector.py}.
	
	\item \textbf{Jetson runtime variant B (TensorRT):} TensorRT engine inference and Flask UI (see also Section~\ref{sec:from_dev_to_dep}).\\
	\noindent Entrypoint: \PATH{Code/LicencePlateDetection/jetson-inference/python/www/flask/dashboard2.py}.
	
	\item \textbf{OCR pipeline:} EasyOCR with preprocessing, normalization, and strict validation.\\
	\noindent Implementation: \PATH{Code/LicencePlateDetection/jetson\_app/ocr\_pipeline.py}.
	
	\item \textbf{Downstream system logic:} Country/format checks, cooldown logic, data logging, and billing logic.\\
	\noindent Module folder: \PATH{Code/LicencePlateDetection/jetson\_app/}.
\end{itemize}

\noindent
In other words, the repository contains both the \textbf{ML model artefacts} and the \textbf{application} that turns detections into an operational demo.
This is important for the delighter comparison: Edge Impulse can replace or simplify parts of the ML workflow, but it does not automatically replace the app.

% --------------------------------------------------------
\section{Edge Impulse overview (what it provides)}
\label{sec:edge_impulse_overview}
\index{Edge Impulse!Studio}
\index{Edge Impulse!CLI}

\noindent
Edge Impulse Studio is designed around projects, datasets, and impulses \cite{EdgeImpulseStudioOverview}.
An \textbf{Impulse} is a packaged pipeline consisting of:
\begin{enumerate}
	\item an \textbf{input block} (e.g., images),
	\item a \textbf{processing block} (feature extraction / preprocessing),
	\item a \textbf{learning block} (training configuration and model family).
\end{enumerate}
\noindent
The Impulse design view is therefore the platform's abstraction for ``the thing that gets deployed'' \cite{EdgeImpulseImpulseDesign}.

\subsection{Data acquisition and labeling}
\index{Edge Impulse!Data acquisition}
\index{Edge Impulse!Labeling}

\noindent
For vision use cases, Edge Impulse supports multiple ingestion routes, such as file upload, API, or CLI tooling \cite{EdgeImpulseDataAcquisition}.
For object detection, it supports bounding box annotations and can import existing labeled datasets in multiple formats \cite{EdgeImpulseUploader,EdgeImpulseObjectDetectionLabelFormats}.
The platform also provides a labeling queue with AI-assisted modes to speed up manual labeling, particularly for bounding boxes \cite{EdgeImpulseLabelingQueue}.

\subsection{Learning blocks and customization options}
\index{Edge Impulse!Learning blocks}
\index{Edge Impulse!BYOM}

\noindent
Edge Impulse provides a set of predefined learning blocks (including object detection blocks), and it also supports advanced customization:
\begin{itemize}
	\item \textbf{Built-in learning blocks:} include object detection (e.g., MobileNet SSD variants, FOMO) and others \cite{EdgeImpulseLearningBlocks}.
	\item \textbf{Custom learning blocks:} can be created (e.g., with PyTorch/Keras) when the predefined options are insufficient \cite{EdgeImpulseLearningBlocks}.
	\item \textbf{Bring Your Own Model (BYOM):} allows importing an externally trained model (e.g., ONNX, SavedModel, LiteRT/TFLite) for profiling, optimization, and deployment \cite{EdgeImpulseBYOM}.
\end{itemize}

\subsection{Deployment artefacts and Linux runners}
\index{Edge Impulse!Deployment}
\index{Edge Impulse!EIM}

\noindent
For Linux devices, Edge Impulse provides a Linux CLI tool suite (including \texttt{edge-impulse-linux-runner}) to run inference on target devices \cite{EdgeImpulseLinuxCLI}.
A central idea is the \textbf{EIM artefact} (Edge Impulse Model), which is compiled for the target architecture and contains the full impulse (processing + learning code) \cite{EdgeImpulseLinuxEIM}.
This changes the integration model: instead of importing a Python model and hand-wiring preprocessing, the developer interacts with an artefact that already contains a standardized pipeline.

\begin{figure}[htbp]
	\centering
	% \usetikzlibrary{positioning,calc}

\begin{tikzpicture}[
	font=\small,
	stage/.style={
		draw, rounded corners, align=center,
		fill=blue!10, text width=42mm,
		inner sep=4pt
	},
	stage2/.style={
		draw, rounded corners, align=center,
		fill=green!12, text width=42mm,
		inner sep=4pt
	},
	arrow/.style={-latex, thick}
	]
	
	% -------- Titles (anchored to first nodes)
	\node[font=\bfseries] (tl) at (-5.1, 2.2) {Our iteration loop};
	\node[font=\bfseries] (tr) at ( 5.1, 2.2) {Edge Impulse iteration loop};
	
	% -------- Left loop (code-first) - relative placement
	\node[stage, below=6mm of tl] (l1) {Collect data\\(camera captures)};
	\node[stage, below=5mm of l1] (l2) {Label \& curate\\(offline tooling)};
	\node[stage, below=5mm of l2] (l3) {Train model\\(\mbox{Ultralytics})};
	\node[stage, below=5mm of l3] (l4) {Export / optimize\\(ONNX / TensorRT)};
	\node[stage, below=5mm of l4] (l5) {Integrate \& deploy\\(Python app)};
	\node[stage, below=5mm of l5] (l6) {Monitor \& debug\\(logs, DB, samples)};
	
	\draw[arrow] (l1) -- (l2);
	\draw[arrow] (l2) -- (l3);
	\draw[arrow] (l3) -- (l4);
	\draw[arrow] (l4) -- (l5);
	\draw[arrow] (l5) -- (l6);
	
	% Clean loop-back on left (symmetric + not too wide)
	\draw[arrow]
	(l6.west) .. controls +(-18mm, 0mm) and +(-18mm, 0mm) .. (l1.west);
	
	% -------- Right loop (platform-first) aligned to left column
	\node[stage2, right=34mm of l1] (r1) {Acquire data\\(Studio / CLI)};
	\node[stage2, below=5mm of r1] (r2) {Label\\(queue + suggestions)};
	\node[stage2, below=5mm of r2] (r3) {Impulse \& training\\(blocks / custom)};
	\node[stage2, below=5mm of r3] (r4) {On-device profiling\\(latency, RAM/ROM)};
	\node[stage2, below=5mm of r4] (r5) {Deployment artifact\\(EIM / SDK)};
	\node[stage2, below=5mm of r5] (r6) {Live validation\\(live classification)};
	
	\draw[arrow] (r1) -- (r2);
	\draw[arrow] (r2) -- (r3);
	\draw[arrow] (r3) -- (r4);
	\draw[arrow] (r4) -- (r5);
	\draw[arrow] (r5) -- (r6);
	
	% Clean loop-back on right (mirror of left)
	\draw[arrow]
	(r6.east) .. controls +(18mm, 0mm) and +(18mm, 0mm) .. (r1.east);
	
	% -------- Shared message: placed relative to bottom of both columns
	\node[
	draw, rounded corners, fill=gray!10, align=left,
	text width=120mm, inner sep=5pt
	] (msg) at ($(l6.south)!0.5!(r6.south) + (0,-16mm)$) {%
		Both loops iterate, but the \textbf{dominant bottleneck} shifts:
		code-first projects spend time on tooling and reproducibility; platform-first
		projects spend time adapting the application around a standardized deployment artefact.};
	
\end{tikzpicture}
	\caption{Iteration loop comparison: code-first workflow vs.\ platform-first workflow.}
	\label{fig:edge_impulse_lifecycle_loop}
\end{figure}

% --------------------------------------------------------
\section{Mapping: Edge Impulse concepts vs.\ our components}
\label{sec:edge_impulse_mapping}
\index{Mapping}

\noindent
Table~\ref{tab:edge_impulse_mapping_table} provides a one-to-one mapping between Edge Impulse workflow concepts and our repository implementation.
This mapping is the basis for a fair comparison: we compare \textbf{equivalent stages}, not unrelated features.

\begin{table}[htbp]
	\centering
	\caption{Mapping between Edge Impulse concepts and our repository components}
	\label{tab:edge_impulse_mapping_table}
	\begin{tabular}{@{}p{0.23\textwidth}p{0.38\textwidth}p{0.33\textwidth}@{}}
		\toprule
		\textbf{System aspect} & \textbf{Our implementation} & \textbf{Edge Impulse equivalent} \\ \midrule
		Data acquisition &
		Manual capture from camera, dataset folder in \FILE{Code/LicencePlateDetection/data/} &
		Data acquisition tab (upload/API/CLI) \cite{EdgeImpulseDataAcquisition,EdgeImpulseUploader} \\ \addlinespace
		Labeling &
		External dataset labels (YOLO format) and offline tools &
		Bounding boxes labeling method + labeling queue \cite{EdgeImpulseLabelingQueue} \\ \addlinespace
		Pipeline definition &
		Scripts + implicit preprocessing choices &
		Impulse (input + processing + learning blocks) \cite{EdgeImpulseImpulseDesign} \\ \addlinespace
		Training &
		Ultralytics training pipeline (YOLO) &
		Learning blocks (built-in, custom blocks, or BYOM) \cite{EdgeImpulseLearningBlocks,EdgeImpulseBYOM} \\ \addlinespace
		Deployment artefact &
		\PYTHON{best.pt} or \PYTHON{best.engine} + custom Python app &
		EIM / SDK artefact built for the target architecture \cite{EdgeImpulseLinuxEIM} \\ \addlinespace
		Application logic &
		OCR, validation, billing, UI (\FILE{Code/LicencePlateDetection/jetson\_app/}) &
		Remains custom (Edge Impulse focuses on the impulse, not the full app) \\ \bottomrule
	\end{tabular}
\end{table}

% --------------------------------------------------------
\section{Comparison dimensions (where Edge Impulse helps, and where it does not)}
\label{sec:edge_impulse_comparison_dimensions}
\index{Comparison}

\subsection{Speed of iteration and reproducibility}
\index{Iteration}

\noindent
In a student project, iteration speed is often determined by hidden overhead:
where datasets live, how labels are managed, how experiments are tracked, and how deployment artefacts are reproduced.
Edge Impulse addresses these problems by providing a centralized Studio workflow for data, labeling, and impulses \cite{EdgeImpulseStudioOverview,EdgeImpulseImpulseDesign}.
In our code-first approach, we achieve transparency and full control, but we also inherit more responsibility:
dataset splits, augmentation decisions, and experiment tracking are enforced by conventions rather than a platform.

\subsection{Model flexibility vs.\ platform constraints}
\index{Flexibility}

\noindent
For object detection, Edge Impulse offers predefined learning blocks and extension paths (custom blocks and BYOM) \cite{EdgeImpulseLearningBlocks,EdgeImpulseBYOM}.
This is powerful, but it also introduces constraints:
for example, imported models have supported formats and defined input/output restrictions that must match the platform's conversion tooling \cite{EdgeImpulseBYOM}.
In our repository, there is no platform-level restriction:
any PyTorch, ONNX, or TensorRT workflow is valid as long as it runs on the Jetson device.

\subsection{Deployment packaging and integration boundary}
\index{Deployment}

\noindent
Edge Impulse changes the integration boundary by delivering a deployment artefact (e.g., EIM) that contains the complete impulse and can be run using standard tools on Linux \cite{EdgeImpulseLinuxCLI,EdgeImpulseLinuxEIM}.
This can reduce the chance of ``it worked on my laptop'' issues because the artefact is compiled for a specific architecture.
However, for our LPR system, the most error-prone part is often not the detector itself, but the \textbf{application glue}:
plate crop selection, OCR preprocessing, post-correction, German format validation, cooldown, persistence, and UI presentation.
These parts remain our responsibility in both workflows.

\begin{figure}[htbp]
	\centering
	% \usetikzlibrary{positioning,calc}

\begin{tikzpicture}[
	font=\small,
	layer/.style={
		draw, rounded corners, fill=blue!8, align=center,
		text width=50mm, inner sep=4pt
	},
	layer2/.style={
		draw, rounded corners, fill=green!10, align=center,
		text width=50mm, inner sep=4pt
	},
	title/.style={font=\bfseries, align=center},
	sep/.style={draw, dashed, gray}
	]
	
	% --- Titles (no absolute coords, anchored)
	\node[title] (tL) {Our deployment stack};
	\node[title, right=60mm of tL] (tR) {Edge Impulse deployment stack};
	
	% --- Common base (aligned automatically)
	\node[layer,  below=6mm of tL] (os1) {Jetson Nano\\JetPack / L4T + CUDA};
	\node[layer2, below=6mm of tR] (os2) {Jetson Nano\\JetPack / L4T + CUDA};
	
	% --- Stack layers (left)
	\node[layer,  below=5mm of os1] (acc1) {Acceleration\\\mbox{TensorRT} engine + \mbox{PyCUDA}};
	\node[layer,  below=5mm of acc1] (ml1)  {Model inference\\\mbox{Ultralytics} or \mbox{TensorRT} bindings};
	\node[layer,  below=5mm of ml1]  (pp1)  {Post-processing\\Crop + OCR + validation};
	\node[layer,  below=5mm of pp1]  (ui1)  {Product surface\\Dashboard + SQLite};
	
	% --- Stack layers (right)
	\node[layer2, below=5mm of os2] (acc2) {Runtime artefact\\EIM / SDK + runner};
	\node[layer2, below=5mm of acc2] (ml2)  {Impulse inference\\(processing + model packaged)};
	\node[layer2, below=5mm of ml2]  (app2) {App integration\\Camera + OCR + business logic};
	\node[layer2, below=5mm of app2] (ui2)  {Product surface\\Your UI + storage};
	
	% --- Separator: exactly spans the stacks (no magic y)
	\draw[sep]
	($(os1.north)!0.5!(os2.north) + (0,3mm)$)
	--
	($(ui1.south)!0.5!(ui2.south) + (0,-3mm)$);
	
\end{tikzpicture}
	\caption{Deployment stack comparison on Jetson Nano: both workflows share the OS/CUDA base, but differ in how the ML pipeline is packaged and consumed by the application.}
	\label{fig:edge_impulse_deploy_stack}
\end{figure}

\subsection{Collaboration and traceability (team scale)}
\index{Collaboration}

\noindent
Edge Impulse includes collaboration-oriented structures such as projects (and, for larger setups, organizations) to centralize datasets and experiments \cite{EdgeImpulseStudioOverview}.
For a group report, this is a delighter because it reduces ``knowledge silo'' risk:
one member does not need to carry the full environment setup to reproduce training results.
In our repository, collaboration is possible via Git and documented scripts, but more discipline is required (tool versions, dataset versioning, and consistent deployment procedures).

% --------------------------------------------------------
\section{Comparison tables (engineering trade-offs)}
\label{sec:edge_impulse_tables}
\index{Tables}

\begin{table}[htbp]
	\centering
	\caption{Qualitative comparison matrix (what we gain and what we trade off)}
	\label{tab:edge_impulse_matrix}
	\begin{tabular}{@{}p{0.30\textwidth}p{0.31\textwidth}p{0.31\textwidth}@{}}
		\toprule
		\textbf{Criterion} & \textbf{Our implementation (code-first)} & \textbf{Edge Impulse (platform-first)} \\ \midrule
		Time-to-first-model &
		Medium (setup + scripts + training) &
		High speed (guided workflow + defaults) \cite{EdgeImpulseStudioOverview} \\ \addlinespace
		Labeling workflow &
		External tooling / dataset-dependent &
		Labeling queue + assisted labeling for bounding boxes \cite{EdgeImpulseLabelingQueue} \\ \addlinespace
		Experiment tracking &
		Script outputs + Git discipline &
		Impulse as a first-class experiment + UI traceability \cite{EdgeImpulseImpulseDesign} \\ \addlinespace
		Deployment packaging &
		Manual exports (\PYTHON{.pt}/\PYTHON{.engine}) + integration &
		EIM/SDK artefacts compiled for target + runner tooling \cite{EdgeImpulseLinuxEIM,EdgeImpulseLinuxCLI} \\ \addlinespace
		OCR \& downstream logic &
		Strong customization (our strength) &
		Still custom (Edge Impulse doesn't replace OCR/business logic) \\ \addlinespace
		Maximum flexibility &
		Very high (any framework/tool) &
		High, but within supported blocks, custom blocks, or BYOM constraints \cite{EdgeImpulseLearningBlocks,EdgeImpulseBYOM} \\ \bottomrule
	\end{tabular}
\end{table}

\subsection{What Edge Impulse does \emph{not} solve for LPR (OCR remains external)}
\label{sec:edge_impulse_ocr_boundary}
\index{OCR}

\noindent
An LPR product is not only a detector; it must output a reliable \textbf{plate string}.
In our system, this reliability is achieved through a chain of OCR-focused decisions:
EU-band masking, multiple thresholding variants, post-correction of common OCR confusions, strict regex validation, and temporal smoothing across frames (\FILE{Code/LicencePlateDetection/jetson\_app/ocr\_pipeline.py}).

\noindent
In an Edge Impulse-first workflow, the object detection part fits naturally into the platform (dataset $\rightarrow$ labels $\rightarrow$ impulse $\rightarrow$ deployment artefact).
However, \textbf{OCR is not inherently solved by object detection}.
For our use case this means:
\begin{itemize}
	\item Edge Impulse can produce a detector that outputs bounding boxes around license plates.
	\item We still need to crop the box, preprocess, run OCR, and validate the result in our application (exactly as we do now).
\end{itemize}

\noindent
This is the key delighter insight: even with a platform that simplifies ML deployment, the \textbf{highest product value} in LPR can sit in the post-processing and domain logic.

\subsection{Practical walkthrough: how we would rebuild the detector with Edge Impulse}
\label{sec:edge_impulse_walkthrough}
\index{Edge Impulse!Walkthrough}

\noindent
To make the comparison actionable, this subsection describes a realistic sequence to migrate only the \textbf{detector lifecycle} to Edge Impulse while keeping OCR and business logic unchanged.
The intent is not to replace our pipeline, but to illustrate \textbf{where effort shifts}.

\begin{enumerate}
	\item \textbf{Export or upload the dataset:} use the Edge Impulse uploader and import bounding boxes in supported formats \cite{EdgeImpulseUploader,EdgeImpulseObjectDetectionLabelFormats}.
	\item \textbf{Validate labeling quality:} correct ambiguous boxes (partial plates, far-distance plates) in the labeling queue \cite{EdgeImpulseLabelingQueue}.
	\item \textbf{Define an impulse:} pick image size and preprocessing blocks suitable for plates (small, high-frequency text) \cite{EdgeImpulseImpulseDesign}.
	\item \textbf{Select a learning block:} start with a predefined object detection block; if accuracy is insufficient, switch to custom blocks or BYOM \cite{EdgeImpulseLearningBlocks,EdgeImpulseBYOM}.
	\item \textbf{Train and evaluate:} compare precision/recall trade-offs (false positives are expensive because OCR will run on every candidate crop).
	\item \textbf{Profile on-device:} use on-device performance reporting to estimate latency and memory on the Jetson class device \cite{EdgeImpulseOnDevicePerformance}.
	\item \textbf{Build deployment artefacts:} export as Linux EIM and test locally on the Jetson Nano using the Linux toolchain \cite{EdgeImpulseLinuxCLI,EdgeImpulseLinuxEIM}.
	\item \textbf{Integrate with our OCR pipeline:} consume bounding boxes from the runner output, then call \PYTHON{PlateOcrPipeline.readPlate()} and downstream logic (\FILE{Code/LicencePlateDetection/jetson\_app/}).
\end{enumerate}

\noindent
The net effect is that Edge Impulse reduces friction in the ML lifecycle (steps 1--7), while our application remains the product layer (step 8).

\subsection{Integration sketch: runner output to our OCR + billing logic}
\label{sec:edge_impulse_integration_sketch}
\index{Edge Impulse!Integration}

\noindent
The integration boundary in a hybrid approach is the detector output.
We already have a well-defined application-side interface: a plate crop (NumPy/OpenCV image) goes into the OCR pipeline and results in a normalized plate string plus confidence.
Therefore, a hybrid integration can be described as: \textbf{replace detector inference call} $\rightarrow$ keep everything downstream.

\begin{lstlisting}[language=Python, caption={Conceptual hybrid integration: Edge Impulse detection output + our OCR pipeline}, label={lst:edge_impulse_hybrid_integration}]
# Pseudo-code sketch (not copy-paste ready)
frame = read_camera_frame()

# 1) Run Edge Impulse runner/EIM to get bounding boxes
bboxes = edge_impulse_detect(frame)  # [(x1, y1, x2, y2, conf), ...]

# 2) For each bbox: crop + OCR + validation (existing code)
for (x1, y1, x2, y2, det_conf) in bboxes:
    plate_crop = frame[y1:y2, x1:x2]
    plate_text, ocr_conf, _ = ocr_pipeline.readPlate(plate_crop)
    plate_text = normalize_and_validate(plate_text)
    if plate_text:
        update_db_and_billing(plate_text, det_conf, ocr_conf)
\end{lstlisting}

\noindent
This sketch highlights why the comparison is meaningful: the \textbf{complexity we already solved} (robust OCR and business logic) remains reusable regardless of the detector tooling.

\subsection{Key takeaways (independent of the chosen implementation path)}
\label{sec:edge_impulse_takeaways}
\index{Takeaways}

\noindent
Even if the final solution remains code-first, Edge Impulse provides reusable ideas:
\begin{itemize}
	\item \textbf{Explicit pipeline definition:} impulse design makes preprocessing and model choices explicit and reviewable \cite{EdgeImpulseImpulseDesign}.
	\item \textbf{Dataset traceability:} centralized labeling workflows reduce annotation drift and ``lost labels'' \cite{EdgeImpulseLabelingQueue}.
	\item \textbf{Early device profiling:} performance reporting encourages designing for latency and memory from day one \cite{EdgeImpulseOnDevicePerformance}.
	\item \textbf{Clear integration boundary:} packaging the impulse into a deployment artefact enforces a stable interface \cite{EdgeImpulseLinuxEIM}.
\end{itemize}

\begin{table}[htbp]
	\centering
	\caption{Proposed evaluation protocol for a fair comparison on Jetson Nano}
	\label{tab:edge_impulse_eval_protocol}
	\begin{tabular}{@{}p{0.26\textwidth}p{0.33\textwidth}p{0.33\textwidth}@{}}
		\toprule
		\textbf{Metric} & \textbf{How to measure (our system)} & \textbf{How to measure (Edge Impulse)} \\ \midrule
		Detection throughput (FPS) &
		Instrument \FILE{Code/LicencePlateDetection/jetson\_app/detector.py} loop with timing and frame counters &
		Run the EIM/runner pipeline and measure inference loop timing \cite{EdgeImpulseLinuxEIM} \\ \addlinespace
		End-to-end latency &
		Time from frame capture $\rightarrow$ OCR result $\rightarrow$ DB event &
		Time from runner input $\rightarrow$ detection output + app-side OCR time \\ \addlinespace
		Memory usage &
		Monitor process RSS + GPU memory (Jetson tools) &
		Use on-device profiling / performance reporting where available \cite{EdgeImpulseBYOM} \\ \addlinespace
		Detection quality &
		mAP/precision/recall on test split (Ultralytics metrics) &
		EI Studio model testing and validation, plus test set analysis \cite{EdgeImpulseStudioOverview} \\ \addlinespace
		OCR quality &
		Character accuracy / plate-string exact match on a labeled OCR test subset &
		Same (OCR remains in the application, independent of EI) \\ \bottomrule
	\end{tabular}
\end{table}

% --------------------------------------------------------
\section{Conclusion}
\label{sec:edge_impulse_conclusion}
\index{Conclusion}

\noindent
From an engineering perspective, Edge Impulse is best seen as a \textbf{workflow accelerator} for the detector lifecycle:
it standardizes ingestion, labeling, impulse design, and deployment artefacts \cite{EdgeImpulseImpulseDesign,EdgeImpulseLinuxEIM}.
For our Jetson Nano LPR system, the delighter insight is that \textbf{the detector is only half of the product}:
the perceived quality comes from robust OCR and application logic under real-world constraints (motion blur, partial occlusion, glare, and repeated reads across frames).

\noindent
Therefore, the best ``hybrid'' strategy for our system is:
\begin{itemize}
	\item Use Edge Impulse to make the detector lifecycle reproducible and faster to iterate (especially for labeling and deployment artefact generation).
	\item Keep OCR, domain validation, billing, persistence, and UI in our codebase where we already have strong customization.
\end{itemize}
